\section{A local refined-decoupling terminal step}
\label{localdec:section}

There is a natural way to insert the refined decoupling theorem of
Guth--Iosevich--Ou--Wang \cite[Theorem 4.2 and Corollary 4.3]{giow}
into the preceding iteration. We prove the resulting estimate here.
The following section shows that its optimized profile cost is already
attainable by the quadratic chain; this particular insertion therefore
does not improve the dimension cutoff.

\begin{lemma}[Weighted local packet estimate]\label{localdec:packet}
Let \(r\gg1\), \(\rho=r^{-1/2}\), and let \(F_\theta\) have Fourier
support in the usual canonical rectangles over a circle of radius
\(r\): tangential width \(O(\rho^{-1})\), normal width \(O(1)\), and
bounded angular overlap. Work in a bounded unit region. Form its usual
smooth packets
\[
 f_T=\chi\eta_{\theta,T}F_\theta,
\]
where \(\chi=1\) on the pin region and the nonnegative transverse
partitions \(\eta_{\theta,T}\) sum to one for each \(\theta\).
Let \(\lambda\) have mass at most one, with density at most \(D\)
after positive smoothing at scale \(r^{-1}\). Retain packets satisfying
\[
 \lambda(CT)\le H\rho,\qquad H\ge1.
\]
The dilation \(CT\) includes the usual small-power packet enlargement
and adjacent \(\rho\)-cubes. For every fixed \(\varepsilon>0\),
\begin{equation}\label{localdec:estimate}
 \int\Big|\sum_{T\ {
m retained}} f_T\Big|^2\,d\lambda
 \lesssim_\varepsilon
 r^\varepsilon H^{2/3}D^{1/3}
 \sum_\theta\int_{Q^+}|F_\theta|^2
 +\operatorname{RapDec}(r).
\end{equation}
As elsewhere, small-power enlargements, polynomial packet counts and
polynomial bounds for the input norms are permitted; their losses are
absorbed by \(r^\varepsilon\).
\end{lemma}

\begin{proof}
Local constancy replaces \(\lambda\) by a positive smooth density
\(w\lesssim D\). Enlarging the test tubes also gives
\(\int_{C'T}w\lesssim H\rho\), up to rapid errors. Pigeonhole the
packet sixth norms, and fix one comparable class of \(W\) packets.
Let \(Y_M\) be a union of \(\rho\)-cubes each meeting between \(M\)
and \(2M\) packets. Weighted incidence counting gives
\[
 M\int_{Y_M}w\lesssim WH\rho.
\]
Indeed every cube incident to a packet lies inside its enlarged test
tube. No equal-mass assumption on these cubes is needed. H\"older and
refined decoupling at exponent six now imply
\begin{align*}
 \int_{Y_M}|f|^2w
 &\le \|f\|_{L^6(Y_M)}^2
          \Big(\int_{Y_M}w^{3/2}\Big)^{2/3}\\
 &\lesssim r^\varepsilon(M/W)^{2/3}
       \sum_T\|f_T\|_6^2\,
       D^{1/3}\Big(\int_{Y_M}w\Big)^{2/3}\\
 &\lesssim r^\varepsilon(H\rho)^{2/3}D^{1/3}
       \sum_T\|f_T\|_6^2.
\end{align*}
The frequency rectangle has area \(O(\rho^{-1})\), so Bernstein gives
\(\|f_T\|_6^2\lesssim\rho^{-2/3}\|f_T\|_2^2\).
For each \(\theta\), bounded overlap of its spatial partition gives
\(\sum_T\|f_T\|_2^2\lesssim\int_{Q^+}|F_\theta|^2\).
The powers of \(\rho\) cancel. Summing the logarithmically many
amplitude and incidence classes proves \eqref{localdec:estimate}.

These are precisely the standard microlocal packets in the cited
corollary. One may use compact smooth spatial cutoffs and their
Schwartz Fourier tails. Alternatively, use nonnegative band-limited
Schwartz transverse partitions and a common band-limited Schwartz
majorant bounded below on the pin region; its lower bound controls
the desired extension there. In this alternative equality to one is
unnecessary. The usual small-power enlargements reduce all tails to
arbitrary inverse powers, while frequency localization, Bernstein and
spatial overlap acquire only small-power losses. Polynomial input norms
and counts allow those errors to be summed. Tiny amplitude classes
are discarded using Bernstein and the same polynomial count.
\end{proof}

Fix one regularized component \(\sigma\) with profile \(f\), and put
\(g(k)=f(k)-k\). For \(R=2^N\) and a parent depth \(m\), set
\[
 b=2^{-m},\quad r=Rb,\quad n=\frac{N+m}{2},\quad
 a=2^{-n}=\sqrt{b/R},\quad \rho=a/b=r^{-1/2}.
\]
Block-grid rounding changes exponents by \(O(T)\). We keep
\(m\le(1-\zeta)N\), with fixed \(\zeta>0\), so the local frequency
is at least \(R^\zeta\). For an active parent \(P\), put
\(\lambda_P=\sigma|_{P^+}/\sigma(P^+)\).
The denominator is at least \(\sigma(P)\). Thus the regular cube
bounds, after rescaling by \(b^{-1}\), give
\begin{equation}\label{localdec:density}
 D_P\lesssim_T R^{C\Gamma}2^{2(N-m)+f(m)-f(N)}
       =R^{C\Gamma}2^{(N-m)+g(m)-g(N)}.
\end{equation}
The old truncated energy bound gives
\[
 J_{\lambda_P}(\rho)\lesssim_T (N+1)R^{C\Gamma}2^{c_g(m,n)},
 \qquad c_g(m,n)=g(m)-\min_{[m,n]}g.
\]
Choose \(H=R^{D_2h+C\Gamma}2^{c_g(m,n)}\).
Lemma~\ref{localdec:packet} then costs
\begin{equation}\label{localdec:terminalcost}
 2^{\mathcal T_g(m)}R^{O(h+\Gamma)+\varepsilon},\qquad
 \mathcal T_g(m)=
 \frac{2c_g(m,(N+m)/2)+(N-m)+g(m)-g(N)}3.
\end{equation}
This uses pin masses at the actual scale \(R^{-1}\), without a
first-cube localization loss.

\paragraph{A single source function for the whole shell.}
Fix the subsequent coarse chain from depth \(m\) to zero. Its inherited
marks are constant on each depth-\(m\) parent \(P\) and are unions of
whole original standard source labels \(S\), of width \(R^{-1/2}\).
Write \(A_{P,S}\in\{0,1\}\) for their survival indicators. Group
these standard labels into nominal ancestors \(\theta\) of width
\(\rho\ge R^{-1/2}\), and form
\[
 F_{P,\theta}=\sum_{S\prec\theta}A_{P,S}F_S.
\]
Decompose these functions into local \emph{pin-side} packets, retaining
whole packets according to the conditional tube bound. Their selected
partition sums
\[
 c_{P,\theta}(y)=\sum_{T\ {
m retained}}\eta_{\theta,T}(y)
 \quad\hbox{satisfy}\quad 0\le c_{P,\theta}\le1.
\]
For \(y\in P\), define
\begin{equation}\label{localdec:goodsource}
 G_{R,y}(x)=\sum_S A_{P,S}c_{P,\theta(S)}(y)
       \chi(x)(\psi_{R,S}\widehat\mu)^\vee(x).
\end{equation}
All angular grids, partitions, test tubes and coefficients are fixed
using \(R\), independently of the circular radius throughout the
annulus. The canonical rectangles have the required dimensions
uniformly for every radius comparable to \(R\), with constants
allowed to depend on its fixed annular ratio. Hence
\eqref{localdec:goodsource} is one source function for the entire shell.
Its coefficients are independent of the distance variable, so for
each fixed pin they commute with the pinned pushforward and circular
identity. Its circular extension on \(P\) is, up to the existing
source-cutoff error, exactly
\(\sum_{\theta,T\ {
m retained}}\eta_{\theta,T}F_{P,\theta}\).

For the discarded first norm, expand only the \emph{original}
unit-length standard source packets. A source--pin pair meets at most
\(R^{O(h)}\) such packets, regardless of how many standard labels are
inside \(\theta\). Their directions lie within
\(O(R^{-1/2+h})\) of the connecting line. Since
\(R^{-1/2}\le\rho\) and \(b\rho=a\), a bad local pin tube implies
the containing heavy tube in that connecting direction. The previous
radial-projection deletion bound therefore applies, with contribution
\[
 \lesssim R^{O(h)}\left(L^{1-q}BR^\Gamma+
       \sum_P\sigma(P^+)\frac{L}{H}J_{\lambda_P}(\rho)\right)
\]
plus the already controlled coarse-chain deletions and rapid errors.
The old choice \(L=R^{D_1h}\), then sufficiently large \(D_2\), makes
this a negative power of \(R\). Short \emph{source} packets would not
justify this argument: after propagation their transverse footprint
is \(\rho\), not \(a\). They have not been used here.

\paragraph{The remaining energy.}
The local selectors disappear from the right side of
\eqref{localdec:estimate}. It leaves precisely the local parent energy
\[
 \sum_{P\ {
m active}}\sigma(P^+)
 \sum_{\theta:\,|\theta|=\rho}
       \int |F_{P,\theta}|^2\,dm_{P^+}.
\]
For the next quadratic inflation, each initial parent label \(\theta\)
is partitioned into descendants of width
\((Rb)^{-1}=\rho^2\le\rho\). The inherited marks were fixed before
the local packet selection, so the old exact inherited identity applies
inside every \(\theta\); no selector remains in this energy.
Continuing the coarse chain and then the original final circle estimate
gives total cost
\begin{equation}\label{localdec:combinedcost}
 \Phi_g(m)+\mathcal T_g(m).
\end{equation}
Here \(\Phi_g(m)\) denotes the cost of the chosen admissible chain from
\(m\) to zero. At \(m=0\), \eqref{localdec:combinedcost} is
\((N-g(N))/3\), which recovers the GIOW shell exponent
\((5-4s)/3\) under the Frostman lower bound. The next result compares
\eqref{localdec:combinedcost} with the fully optimized quadratic chain.

\paragraph{The weighted family below exponent six.}
The same packet-compatible construction allows the weighted theorem of
Du--Ou--Ren--Zhang \cite[Theorem 1.1(b)]{dorz} to be used without
throwing away its natural-scale gain. In unit parent coordinates,
suppose the smoothed pin density satisfies \(w\le A\), and every
\(\rho\)-square has pin mass at most \(v\). Put
\[
 Q=\frac{rv}{A}\le1.
\]
In large coordinates \(X=rx\), the weight \(w(X/r)/A\) lies in
\([0,1]\) and has mass at most \(r^2v/A\) on each
\(r^{1/2}\)-square. Taking
\(\alpha=2\log(r^2v/A)/\log r\le2\) in that theorem gives the
squared weighted gain \(Q^{2/p-1/3}\), for \(2\le p\le6\).
For one comparable packet class of cardinality \(W\), H\"older with
respect to \(w\,dx\), followed by that theorem, yields
\[
 \int_Y|f|^2w
 \lesssim r^\varepsilon A^{2/p}Q^{2/p-1/3}
 \left(\frac{M\int_Yw}{W}\right)^{1-2/p}
 \sum_T\|f_T\|_p^2.
\]
The same incidence bound and the packet Bernstein estimate
\(\|f_T\|_p^2\lesssim\rho^{-(1-2/p)}\|f_T\|_2^2\)
therefore give multiplier
\begin{equation}\label{localdec:weightedmultiplier}
 A^{2/p}Q^{2/p-1/3}H^{1-2/p}
\end{equation}
for the local parent square function. All counts and small-power losses
are handled as in Lemma~\ref{localdec:packet}.

Writing \(L=N-m\), the leading profile exponents are
\[
 I=L+g(m)-g(N),\qquad
 J=-L/2+g(N)-g(n),\qquad c=c_g(m,n),
\]
\[
 A=2^I,\quad v=2^{-L/2+g(m)-g(n)},\quad Q=2^J,\quad H=2^c.
\]
Here \(J\le0\) by the Lipschitz inequality between \(n\) and \(N\).
The block and regularization errors only add the usual arbitrarily
small losses; enlarging \(A\) slightly also enforces \(Q\le1\)
when those errors are present. Thus the weighted terminal cost is
\begin{equation}\label{localdec:weightedcost}
 \mathcal T_{g,p}(m)=\frac2p I+
        \left(\frac2p-\frac13\right)J+
        \left(1-\frac2p\right)c.
\end{equation}
At \(p=6\) this is \eqref{localdec:terminalcost}. The comparison with
other exponents and with optimized quadratic chains is purely a
profile calculation, proved next. These statements concern the
specific terminal construction just described, not all possible uses
of refined or weighted decoupling.
